\section{Discussion and Conclusion}\label{sec:discussion}

\paragraph{Implications and novelty.} Our central finding, that metric choice alone produces a 14--61 layer divergence and qualitatively reverses hidden-state advantage conclusions, means safety probing studies should report results under multiple layer selection criteria. The cosine compression mechanism (\cref{app:cosine_volatility}) provides a geometric explanation: shallow-layer representations compress directional information, inflating crossing-rate metrics where probes operate near the noise floor. This mechanism is architecture-general (confirmed on Gemma-2; \cref{app:gemma_cross_arch}), though tested only on standard attention models. While metric sensitivity is well-documented~\citep{belinkov2022probing,voita2020information}, our contribution is demonstrating a \emph{cascade} specific to temporal safety probing: metric $\to$ layer $\to$ conclusion, where the severity (14--61 layer gap, sign reversal) exceeds what prior work documents at a fixed layer.\looseness=-1

\paragraph{Practical recommendations.} We propose three levels of defense against metric-induced layer selection bias (\cref{app:practical_mitigation}). \emph{Level~1 (minimal)}: report layer selections under both threshold-based and frequency-based metric families; divergence exceeding 20\% of model depth indicates metric-induced bias. \emph{Level~2 (recommended)}: compute all five metrics and flag when ${\geq}3$ select layers differing by ${>}20\%$ of depth. \emph{Level~3 (robust)}: a three-layer ensemble averaging predictions from metric-selected layers eliminates pre-selection entirely ($p \leq 0.002$; \cref{tab:ensemble}), but requires cross-validated confirmation.\looseness=-1

Table~\ref{tab:consolidated} consolidates per-model findings across all three contributions.

\begin{table}[t]
\centering
\small
\resizebox{\columnwidth}{!}{%
\begin{tabular}{lccccl}
\toprule
\textbf{Model} & \textbf{Layer gap} & $\Delta$\textbf{Prec (384d)} & $\Delta$\textbf{Prec (1024d)} & \textbf{Pos.\ resid.} & \textbf{Evidence} \\
\midrule
R1-8B & L0$\to$L14 & +20.1pp \checkmark & +5.8pp $\times$ & Lin \checkmark\ Quad \checkmark & Cross-val. \\
OT-7B & L0$\to$L16 & +21.9pp \checkmark & +18.1pp \checkmark & Lin \checkmark\ Quad \checkmark & Cross-val. \\
QwQ-32B$^\dagger$ & L2$\to$L63 & +19.2pp \checkmark & +9.4pp $\times$ & Lin \checkmark\ Quad $\times$ & Suggestive \\
R1-32B$^\dagger$ & L10$\to$L63 & +15.4pp \checkmark & +11.7pp $\times$ & $\times$\ $\times$ & Suggestive \\
\bottomrule
\end{tabular}}
\caption{Consolidated findings across contributions. Layer gap: crossing-rate composite vs.\ threshold-based selection (C1). $\Delta$Prec: HS$-$text precision advantage at threshold-based layers; \checkmark/{\scriptsize$\times$} = survives/fails Holm--Bonferroni (C2). Pos.\ resid.: HS advantage survives linear/quadratic position residualization. $^\dagger$Sample-limited ($n \leq 20$).}
\label{tab:consolidated}
\end{table}

\paragraph{Boundary conditions.} The hidden-state advantage is confirmed at 7--8B (+20--22pp, cross-validated with 13--25\% nested CV shrinkage; \cref{app:nested_cv}) but conditional on several factors. It diminishes from 4/4 to 1/4 significant models with stronger encoders (384d $\to$ 1024d), suggesting that as text encoders improve, the marginal value of hidden-state access decreases. Position partially mediates the advantage: under linear residualization, 3/4 models retain significant non-positional signal (+3.7 to +6.9pp after correction), but R1-32B's residualized gap (+1.9pp) is not significant; quadratic residualization absorbs the advantage in 2/4 models (QwQ-32B and R1-32B fully; R1-8B and OT-7B retain significance; \cref{tab:position_quadratic}). At 32B ($n \leq 20$), effects below 8.6--9.5pp cannot be reliably distinguished from noise at 80\% power (\cref{app:power_analysis}); replication with $n \geq 50$ traces is needed. HarmThoughts is currently the only dataset with step-level CP annotations; AdvBench confirms the geometric mechanism generalizes but does not validate the hidden-state advantage. Ensemble layer selection (\cref{tab:ensemble}) mitigates the bias but false positive rates (11--50\%) preclude deployment.\looseness=-1

\paragraph{Future work.} Causal interventions on position encoding (e.g., shuffled or ablated position embeddings) would disentangle the content-level safety signal from temporal artifacts. Extending to online detection with streaming probes and calibrated confidence thresholds would move from retrospective characterization toward deployment.\looseness=-1
